\documentclass[11pt]{book}
\usepackage{fontspec}
\setmainfont{Latin Modern Roman}
\setsansfont{Latin Modern Sans}
\usepackage{amsmath,amssymb,mathtools}
\usepackage{unicode-math}
\setmathfont{Latin Modern Math}
\usepackage{xurl}
\usepackage[hidelinks]{hyperref}
\newcommand{\sref}[2]{\hyperlink{source:#1:#2}{[S#2]}}
\newcommand{\cataloguescope}[1]{\paragraph{Mathematical question.}#1}
\begin{document}
\section{Research Problems in Function Theory — Problem 9.9}\label{statement}

\subsection{Definitions and mathematical statement}

\paragraph{Shared notation.}$\mathbb N=\{1,2,\ldots\}$, and $\mathbb Z,\mathbb Q,\mathbb R,\mathbb C$ denote the integers, rationals, reals and complex numbers. Any other conventions are specified locally.


Let $(p_n)_{n\ge1}$ be positive integers with $\sum_{n\ge1}p_n^{-1}<\infty$. For a Jordan arc $\gamma\subset\mathbb C$, the space $C(\gamma)$ has its uniform norm; the powers span it when finite complex linear combinations of the restrictions $z^{p_n}|_\gamma$ are uniformly dense. No constant monomial is added unless it occurs in the specified family, and all $p_n$ are positive.
\paragraph{Statement recorded in the supplied source.}
Does reciprocal summability imply that the powers fail to span $C(\gamma)$ for every planar Jordan arc $\gamma$? \sref{statement}{2}
\subsection{Short English statement}
Must reciprocal-summable positive powers fail to be dense in the continuous functions on every planar Jordan arc?
\subsection{Sources}
\begin{description}
\item[\hypertarget{source:statement:1}{S1}] ULAM.AI and the UnsolvedMath Contributors, \emph{UnsolvedMath: A Curated Collection of Open Mathematics Problems}, catalogue entry. CC BY 4.0. \url{https://huggingface.co/datasets/ulamai/UnsolvedMath}.
\item[\hypertarget{source:statement:2}{S2}] Walter K. Hayman and Edward F. Lingham, \emph{Research Problems in Function Theory} (2018), arXiv:1809.07200. Problem 9.9, complete original question, all following updates, and the relevant chapter definitions. \url{https://arxiv.org/abs/1809.07200}.
\end{description}
\label{end:statement}
\end{document}
